\chapter{Torticollis wry neck}
\index{Torticollis wry neck}

\section*{Definition and Overview}
Torticollis, derived from the Latin words \textit{tortus} (meaning twisted) and \textit{collum} (meaning neck), is a clinical sign characterized by an abnormal, asymmetrical tilt, rotation, or flexion of the head and neck. Commonly referred to as "wry neck," it is not a distinct disease entity but rather a manifestation of a wide variety of underlying musculoskeletal, neurological, infectious, ocular, or traumatic etiologies. The hallmark of the condition is the involuntary contracture or spasm of the cervical muscles, most notably the sternocleidomastoid (SCM) muscle, which results in the lateral tilting of the head toward the affected side and rotation of the chin toward the opposite side.

Torticollis can be broadly classified based on its onset into congenital and acquired forms. Congenital muscular torticollis (CMT) is the most common presentation, typically diagnosed within the first few weeks to months of life. It arises from localized changes within the SCM muscle during fetal development or delivery, resulting in a fibrotic, shortened muscle band. Acquired torticollis, conversely, develops later in life, spanning from infancy through adulthood. Acquired forms can arise suddenly, as in the case of acute idiopathic wry neck, or manifest as a chronic, progressive condition, such as cervical dystonia (spasmodic torticollis).

Understanding the scope of torticollis is clinically significant. While many cases of acute acquired torticollis are benign, self-limiting, and resolve within days with conservative care, other presentations can signal life-threatening conditions. For instance, acute-onset torticollis in a pediatric patient may indicate a deep neck infection (e.g., retropharyngeal abscess), atlantoaxial subluxation, or a posterior fossa tumor. In adults, chronic cervical dystonia can cause severe, disabling pain and profound psychosocial distress. Therefore, a comprehensive understanding of its pathophysiology, clinical features, and differential diagnosis is essential for medical professionals to rule out serious underlying pathology and implement appropriate, timely interventions.

\section*{Epidemiology}
The epidemiology of torticollis varies significantly depending on the classification, clinical subtype, and patient age group. Congenital muscular torticollis (CMT) is one of the most common musculoskeletal anomalies in newborns, third only to developmental dysplasia of the hip (DDH) and congenital talipes equinovarus (clubfoot). Studies report the incidence of CMT to range between 0.3\% and 2.0\% of all live births. There is a documented association between CMT and primiparous mothers, fetal breech presentation, prolonged or difficult labor, and assisted deliveries involving forceps or vacuum extraction. Some epidemiological data suggest a slight male predominance, though findings remain variable across different geographic and demographic cohorts. Up to 10\% to 20\% of infants diagnosed with CMT also exhibit coexisting developmental dysplasia of the hip, highlighting the necessity of clinical screening for associated orthopedic conditions.

Acquired torticollis has a broader, more heterogeneous distribution. Acute wry neck, often referred to as acute positional or idiopathic torticollis, is highly prevalent in children, adolescents, and young adults. Because many individuals manage mild, transient episodes at home without seeking medical attention, the true community incidence is difficult to determine, but it remains a frequent complaint in primary care and emergency departments.

Spasmodic torticollis, or cervical dystonia, is the most common form of focal dystonia. It typically presents in adulthood, with a peak age of onset between 30 and 50 years. Epidemiological studies estimate the prevalence of cervical dystonia to be between 9 and 30 cases per 100,000 individuals in Western populations. Unlike CMT, cervical dystonia exhibits a prominent female-to-male ratio of approximately 2:1. The disease can be sporadic or, in up to 10\% to 15\% of cases, familial, demonstrating an autosomal dominant inheritance pattern with reduced penetrance.

\section*{Aetiology and Pathophysiology}
The etiologies of torticollis are diverse and complex, spanning musculoskeletal, neurological, infectious, ocular, and pharmacological domains. The pathophysiology depends entirely on the specific underlying etiology:

\begin{itemize}
    \item \textbf{Congenital Muscular Torticollis (CMT):} The primary pathophysiology involves localized fibrosis, shortening, and contracture of the sternocleidomastoid (SCM) muscle. The exact cause remains debated, but the leading theory involves intrauterine packaging or fetal malpositioning, particularly in breech presentations, which causes persistent lateral flexion of the neck. This positioning can lead to localized compartment syndrome, venous congestion, and ischemia within the SCM muscle belly. Alternatively, birth trauma resulting from excessive traction on the neck during delivery can cause microvascular injury and intramuscular hematoma. Over time, the damaged muscle fibers undergo necrosis, which is replaced by dense fibrotic collagen tissue. This process frequently manifests as a palpable, non-tender fibrous mass, or "pseudotumor," within the SCM muscle belly during the first few weeks of life, which eventually resolves into a tight, inelastic band.
    \item \textbf{Cervical Dystonia (Spasmodic Torticollis):} This is a chronic, neurological movement disorder characterized by involuntary, sustained, or intermittent muscle contractions of the neck. The pathophysiology is primarily attributed to dysfunction within the basal ganglia-thalamocortical circuits, which are responsible for motor planning, execution, and execution inhibition. There is a disruption in reciprocal inhibition, leading to the simultaneous contraction of agonist and antagonist muscle groups. Neurochemical studies indicate alterations in neurotransmitter pathways, particularly dopaminergic dysfunction and cholinergic hyperactivity. A genetic component is also recognized; mutations in genes such as \textit{TOR1A} (linked to early-onset generalized dystonia DYT1), \textit{THAP1} (DYT6), and \textit{GNAL} (DYT25) have been implicated in familial forms of the disease.
    \item \textbf{Infectious and Inflammatory Pathways (Grisel's Syndrome):} This is a non-traumatic subluxation of the atlantoaxial (C1-C2) joint, occurring primarily in children. The pathophysiology involves the hematogenous spread of infection or inflammation from adjacent structures, such as the pharynx, tonsils, adenoids, or middle ear, via the pharyngovertebral venous plexus. The localized inflammatory response leads to hyperemic decalcification of the anterior arch of the atlas and laxity of the transverse and alar ligaments. The weakened ligaments can no longer stabilize the atlantoaxial joint, resulting in subluxation and subsequent reactive muscle spasm of the SCM and surrounding paraspinal muscles, locking the head in a rotated posture.
    \item \textbf{Ocular Torticollis:} This is a compensatory, non-musculoskeletal head tilt adopted by patients to optimize binocular vision, maintain alignment of the visual axes, or minimize diplopia (double vision). It most commonly occurs in the setting of congenital or acquired fourth cranial nerve (trochlear nerve) palsy, which paralyzes the superior oblique muscle. Because the superior oblique muscle is responsible for intorsion, depression, and abduction of the eye, a palsy causes the affected eye to extort and elevate. The patient tilts their head toward the opposite shoulder to align the eyes and prevent double vision. Other causes include congenital nystagmus or strabismus, where the head tilt brings the eyes into a "null zone" of minimal nystagmus or optimal visual acuity.
    \item \textbf{Traumatic Musculoskeletal Injury:} Direct trauma to the cervical spine, such as whiplash injuries, minor sprains of the cervical facet joints, or microtrauma to the paraspinal muscles, can trigger a protective reflex muscle spasm. The nervous system restricts motion to prevent further injury to the damaged tissues, locking the cervical spine in a painful, asymmetrical posture.
    \item \textbf{Drug-Induced Acute Dystonic Reactions:} This is a pharmacological etiology resulting from the administration of dopamine receptor antagonists. The most common offending agents are first-generation antipsychotics (e.g., haloperidol, chlorpromazine, fluphenazine) and certain antiemetic medications (e.g., metoclopramide, prochlorperazine). The pathophysiology involves the sudden, acute blockade of post-synaptic D2 dopamine receptors in the nigrostriatal pathway. This disruption causes an acute imbalance between dopaminergic (inhibitory) and cholinergic (excitatory) neurotransmission, triggering severe, involuntary, and painful spasms of the neck muscles (retrocollis or torticollis), jaw (trismus), tongue, and facial muscles.
\end{itemize}

\section*{Clinical Features}
The clinical presentation of torticollis varies based on whether the condition is congenital or acquired, acute or chronic, and its primary underlying cause. However, certain physical findings are common to many forms:

\begin{itemize}
    \item \textbf{Asymmetrical Head Posture:} The classic presentation is a lateral tilt of the head toward the affected SCM muscle, accompanied by rotation of the face and chin toward the contralateral (opposite) shoulder. The shoulder on the affected side may also be elevated. In less common variants, such as retrocollis or antecollis, the head is pulled directly backward or forward, respectively.
    \item \textbf{Restricted Range of Motion:} Patients exhibit a significant limitation in both active and passive range of motion of the cervical spine. Specifically, there is restricted passive rotation of the head toward the ipsilateral (affected) side and restricted lateral flexion toward the contralateral side.
    \item \textbf{Sternocleidomastoid Muscle Pathology:} In infants with congenital muscular torticollis, a palpable, firm, olive-shaped mass (typically 1 to 3 cm in diameter) may be detected within the lower half of the SCM muscle between 2 and 8 weeks of age. This mass is non-tender and does not exhibit signs of acute inflammation (such as warmth or erythema). Over several months, this mass typically regresses, leaving a tight, inelastic fibrous band that limits muscle excursion.
    \item \textbf{Craniofacial Asymmetry (Plagiocephaly):} If CMT is left untreated in infants, persistent unilateral positioning leads to deformational plagiocephaly. This is characterized by flattening of the parieto-occipital skull on the contralateral side, anterior displacement of the ipsilateral ear, and hemifacial microsomia (flattening of the forehead, cheek, and jaw on the affected side).
    \item \textbf{Pain and Muscle Spasms:} In acute acquired torticollis, patients present with a sudden onset of sharp, severe, unilateral neck pain that is exacerbated by any attempt to move the head. The affected muscles are exquisitely tender to palpation and feel rigid or "locked." Conversely, infants with CMT rarely experience pain. In adults with chronic cervical dystonia, the pain is described as a deep, aching, or burning sensation in the neck and shoulders, resulting from constant, involuntary muscle contractions.
    \item \textbf{Neurological and Sensory Features:} Patients with cervical dystonia may exhibit involuntary head tremors or jerking movements (clonic spasms). A classic feature of cervical dystonia is the \textit{geste antagoniste} (sensory trick), where a light touch to the chin, cheek, or back of the head temporarily reduces the severity of the dystonic muscle contraction, allowing the patient to align their head.
    \item \textbf{Systemic and Infectious Signs:} In cases of infectious origin (e.g., Grisel's syndrome, retropharyngeal abscess), the head tilt is accompanied by systemic symptoms such as high fever, sore throat, painful swallowing (odynophagia), difficulty swallowing (dysphagia), drooling, muffled voice ("hot potato voice"), and tender cervical lymphadenopathy.
\end{itemize}

\section*{Differential Diagnosis}
A broad range of conditions can cause or mimic torticollis. A careful diagnostic approach is necessary to distinguish benign musculoskeletal conditions from serious neurological, infectious, or neoplastic diseases:

\begin{itemize}
    \item \textbf{Congenital Muscular Torticollis (CMT):} The primary diagnosis in infants presenting with early-onset head tilt and SCM tightness. It must be differentiated from other congenital anomalies.
    \item \textbf{Atlantoaxial Rotatory Subluxation (AARS):} A mechanical displacement of the C1 and C2 vertebrae. It can occur secondary to minor trauma or secondary to pharyngeal inflammation (Grisel's syndrome). AARS presents with a persistent "cock-robin" posture (head tilted to one side, rotated to the opposite side, and slightly flexed).
    \item \textbf{Sandifer Syndrome:} A pediatric condition that mimics neurological torticollis. It is characterized by paroxysmal dystonic postures of the neck and back (opisthotonus) associated with gastroesophageal reflux disease (GERD) or hiatal hernia. The movements are self-limiting and occur almost exclusively during or shortly after feeding.
    \item \textbf{Klippel-Feil Syndrome:} A congenital skeletal anomaly characterized by the improper segmentation and fusion of two or more cervical vertebrae. The clinical triad includes a short neck, a low posterior hairline, and a severely restricted range of neck motion.
    \item \textbf{Posterior Fossa and Spinal Cord Neoplasms:} Tumors of the cerebellum, brainstem, or cervical spinal cord (e.g., astrocytoma, ependymoma, medulloblastoma) can cause torticollis. This is due to direct irritation of the accessory nerve or the vestibular pathways, or as a compensatory response to raise intracranial pressure. It is typically accompanied by progressive headache, vomiting, ataxia, or focal neurological deficits.
    \item \textbf{Ocular Motility Deficits (Ocular Torticollis):} Abnormal head posture adopted to compensate for strabismus, congenital nystagmus, or superior oblique muscle palsy (fourth cranial nerve palsy). In these cases, there is no intrinsic restriction of the cervical muscles, and the head tilt disappears when one eye is patched or when the patient is asleep.
    \item \textbf{Deep Neck Space Infections:} Retropharyngeal, parapharyngeal, or peritonsillar abscesses, and acute cervical lymphadenitis. These infections cause intense, localized tissue inflammation, leading to reactive spasm of the adjacent prevertebral and paraspinal muscles.
    \item \textbf{Drug-Induced Acute Dystonic Reaction:} Extrapyramidal side effects from dopamine receptor antagonists. The sudden onset of painful neck spasms, facial grimacing, and oculogyric crisis (involuntary upward deviation of the eyes) is diagnostic, especially with a history of recent medication initiation or dosage increase.
    \item \textbf{Benign Paroxysmal Torticollis of Infancy:} A rare, self-limiting disorder characterized by recurrent, transient episodes of head tilting toward one side. The episodes last from a few hours to several days and may be accompanied by vomiting, pallor, irritability, and ataxia. The condition resolves spontaneously, usually by the age of 2 to 3 years.
\end{itemize}

\section*{When to Seek Medical Care}
While acute positional torticollis often resolves with conservative measures, certain signs and symptoms require prompt medical evaluation. Parents of infants showing any persistent head tilt or asymmetry should consult their pediatrician.

\noindent\textbf{IMPORTANT WARNING: Immediate emergency medical attention must be sought if torticollis is accompanied by any of the following "red flag" symptoms:}

\begin{itemize}
    \item \textbf{Airway Compromise:} Difficulty breathing, stridor (a high-pitched gasping sound during inhalation), drooling, or an inability to swallow secretions, which may indicate a rapidly progressing deep neck infection (e.g., retropharyngeal abscess).
    \item \textbf{Severe Post-Traumatic Onset:} Sudden onset of torticollis following a motor vehicle collision, fall, or sports-related injury, which raises concern for a cervical spine fracture, dislocation, or atlantoaxial instability.
    \item \textbf{Focal Neurological Deficits:} Weakness, numbness, or paresthesia (tingling) in the arms or legs, changes in gait, or loss of bowel or bladder control, suggesting spinal cord or nerve root compression.
    \item \textbf{Systemic Infectious Symptoms:} High fever, severe headache, photophobia (sensitivity to light), altered mental status, or significant neck rigidity when flexing the chin toward the chest, which are highly suspicious for meningitis.
\end{itemize}

For emergency assistance in life-threatening scenarios, immediately call the appropriate emergency number. In many countries, call local emergency services. In the United States and Canada, call 911. In the United Kingdom, call 999. In Europe, call 112.

\section*{Diagnosis and Investigations}
The diagnosis of torticollis is primarily clinical, relying on a thorough medical history and a systematic physical examination. Investigations are reserved to confirm the underlying etiology, assess the severity of muscle involvement, and rule out structural or neurological pathology.

\begin{itemize}
    \item \textbf{Clinical Assessment:}
    \begin{itemize}
        \item \textit{History:} In infants, detailed history of maternal pregnancy (breech presentation, multiple gestations), delivery method (forceps, vacuum, Caesarean section), and when the head tilt was first observed. In adults, history of trauma, occupational neck strain, family history of movement disorders, and a thorough review of current medications.
        \item \textit{Physical Examination:} Assessment of active and passive cervical range of motion (rotation, lateral flexion, extension). Palpation of the sternocleidomastoid and upper trapezius muscles to detect tenderness, localized spasm, or fibrous bands. A thorough neurological exam assessing cranial nerves, motor strength, sensory function, and deep tendon reflexes is mandatory.
        \item \textit{Ophthalmic Screen:} Evaluation of extraocular movements, pupillary reflexes, and visual acuity to rule out ocular causes of head tilt.
    \end{itemize}
    \item \textbf{Imaging Modalities:}
    \begin{itemize}
        \item \textit{Ultrasonography:} This is the gold standard imaging choice for infants suspected of congenital muscular torticollis. It is a non-invasive, radiation-free modality that can visualize the internal architecture of the SCM muscle. Fibrosed SCM tissue typically appears as a localized hyperechoic mass or fusiform muscle thickening.
        \item \textit{Cervical Spine Radiography (X-rays):} Standard anteroposterior (AP), lateral, and open-mouth (odontoid) views are indicated in patients with a history of trauma, acute onset without muscle spasm, or suspected bony anomalies. X-rays are useful to identify cervical vertebrae fusion (Klippel-Feil syndrome), hemivertebrae, osteomyelitis, or atlantoaxial subluxation.
        \item \textit{Computed Tomography (CT) Scanning:} CT provides detailed resolution of bony structures and is the diagnostic tool of choice for demonstrating atlantoaxial rotatory subluxation (AARS). Contrast-enhanced CT is also invaluable for diagnosing deep neck space infections like retropharyngeal abscesses.
        \item \textit{Magnetic Resonance Imaging (MRI):} Indicated if there are signs of central nervous system involvement, progressive neurological deficits, or when a posterior fossa/spinal cord tumor is suspected. MRI provides superior soft-tissue resolution of the brainstem, cerebellum, and spinal cord.
    \end{itemize}
    \item \textbf{Laboratory Investigations:} Complete blood count (CBC), erythrocyte sedimentation rate (ESR), and C-reactive protein (CRP) are indicated if an infectious or inflammatory process is suspected based on fever, pharyngitis, or localized lymphadenopathy.
\end{itemize}

\section*{Management}
The management of torticollis is highly individualized, tailored to the specific underlying cause, the age of the patient, and the severity and duration of the symptoms. Treatment options span conservative, pharmacological, and surgical domains.

\begin{itemize}
    \item \textbf{Conservative and Physical Therapy:}
    \begin{itemize}
        \item \textit{Congenital Muscular Torticollis (CMT):} Early intervention is the key to successful conservative management. Physical therapy consists of gentle, passive stretching of the tight SCM muscle, active neck rotation exercises, and developmental positioning. Parents are instructed on positioning techniques, such as placing the infant's crib or changing table in a way that forces them to turn their head away from the affected side to see parents or toys. Supervised prone positioning ("tummy time") during waking hours is critical to strengthen neck musculature and prevent positional plagiocephaly. Over 90\% of infants resolve completely with physical therapy when initiated before 3 months of age.
        \item \textit{Acute Acquired Torticollis:} Treatment includes local application of moist heat, gentle passive stretching within pain limits, and temporary soft cervical collar support (limited to a few days to avoid muscle atrophy). Gentle manual cervical traction and massage can also relieve acute spasm.
    \end{itemize}
    \item \textbf{Pharmacological Therapy:}
    \begin{itemize}
        \item \textit{Analgesics and Anti-inflammatories:} Nonsteroidal anti-inflammatory drugs (NSAIDs) such as ibuprofen, or simple analgesics like paracetamol, are first-line agents to control pain in acute musculoskeletal wry neck.
        \item \textit{Muscle Relaxants:} Short courses of muscle relaxants (e.g., diazepam, baclofen, cyclobenzaprine) can be prescribed to reduce painful muscle spasms.
        \item \textit{Anticholinergic Agents:} Medications like trihexyphenidyl may be used in adults with cervical dystonia, though their efficacy is often limited by systemic anticholinergic side effects (e.g., dry mouth, blurred vision, urinary retention).
        \item \textit{Antidotes for Drug-Induced Reactions:} Acute dystonic reactions from dopamine antagonists represent a medical emergency. They are treated with parenteral administration of anticholinergic agents like benztropine (1 to 2 mg IV or IM) or antihistamines with strong anticholinergic properties like diphenhydramine (50 mg IV or IM). Symptoms typically resolve within minutes of administration.
    \end{itemize}
    \item \textbf{Botulinum Toxin Injections:} Botulinum toxin type A (BoNT-A) is the gold standard, first-line medical therapy for chronic cervical dystonia and is also used in cases of CMT refractory to 6 months of stretching. The toxin is injected directly into the affected muscles (most commonly the SCM, splenius capitis, and upper trapezius), often guided by electromyography (EMG) or ultrasound. BoNT-A inhibits the release of acetylcholine at the neuromuscular junction, causing temporary muscle relaxation. The clinical benefit typically manifests within 1 to 2 weeks and lasts for approximately 3 to 4 months, requiring repeated injections.
    \item \textbf{Surgical Intervention:}
    \begin{itemize}
        \item \textit{Surgical Release for CMT:} Indicated for infants who do not respond to conservative physical therapy after 12 to 18 months, or in children presenting after 1 year of age with severe SCM contracture and progressive facial asymmetry. Surgical options include a distal unipolar release, bipolar release (releasing both the mastoid and clavicular/sternal attachments of the SCM), or Z-plasty lengthening of the muscle. Postoperative physical therapy and bracing are crucial to maintain alignment.
        \item \textit{Interventions for Cervical Dystonia:} In severe, medically refractory cervical dystonia, surgical options include selective peripheral denervation (excision of the distal motor branches of the accessory nerve and posterior primary rami of C1 to C3/C4). Alternatively, Deep Brain Stimulation (DBS) targeting the globus pallidus internus (GPi) has emerged as an effective neurosurgical option for restoring motor control.
    \end{itemize}
\end{itemize}

\section*{Complications}
The complications of torticollis depend on the patient's age at onset, the underlying cause, and whether prompt and appropriate treatment was initiated.

\begin{itemize}
    \item \textbf{Complications of Untreated CMT:}
    \begin{itemize}
        \item \textit{Deformational Plagiocephaly:} Persistent pressure on one side of the soft infant skull causes severe, permanent flattening of the occipital bone and asymmetry of the forehead, ears, and orbits.
        \item \textit{Facial Hemihypoplasia:} Lack of muscle excursion and persistent tilt restrict blood flow and normal mechanical forces on the facial bones, leading to underdevelopment of the cheek, jaw, and eye socket on the affected side.
        \item \textit{Compensatory Scoliosis:} The child develops abnormal lateral curvature of the thoracic and cervical spine as the body attempts to realign the eyes horizontally with the horizon.
        \item \textit{Motor Developmental Delays:} The restricted range of motion limits visual tracking and asymmetrical weight-bearing, which can delay motor milestones such as rolling, crawling, and sitting.
    \end{itemize}
    \item \textbf{Complications of Untreated Cervical Dystonia:}
    \begin{itemize}
        \item \textit{Cervical Spondylosis and Osteoarthritis:} Chronic, repetitive, and forceful muscle contractions exert asymmetric mechanical stress on the cervical vertebrae, leading to premature degeneration of intervertebral discs, facet joint hypertrophy, and osteophyte formation.
        \item \textit{Cervical Radiculopathy and Myelopathy:} Degenerative changes or disc herniations secondary to dystonic movements can compress cervical nerve roots or the spinal cord, leading to chronic radicular pain, sensory loss, motor weakness, or spastic gait.
        \item \textit{Chronic Pain and Disability:} Persistent muscle spasms cause intractable, exhausting pain, leading to significant limitations in daily activities, employment, and social functioning, often culminating in secondary depression and anxiety.
    \end{itemize}
    \item \textbf{Treatment-Related Complications:}
    \begin{itemize}
        \item \textit{Surgical Risks:} These include injury to the accessory nerve or facial nerve, damage to the jugular veins or carotid arteries, wound infection, cosmetically displeasing scars, or failure to correct the deformity.
        \item \textit{Botulinum Toxin Adverse Effects:} Diffusion of the toxin into adjacent pharyngeal muscles can cause transient dysphagia (difficulty swallowing), dysphonia (hoarseness), or severe neck muscle weakness, causing the head to tilt in the opposite direction.
    \end{itemize}
\end{itemize}

\section*{Prognosis and Outcomes}
The prognosis for individuals diagnosed with torticollis is generally highly favorable, particularly when the condition is identified early and managed with appropriate therapeutic strategies.

In congenital muscular torticollis (CMT), the timing of treatment initiation is the most critical prognostic factor. When passive stretching and positioning therapies are started within the first three months of life, full clinical resolution is achieved in over 90\% to 95\% of cases, with normal range of motion and no residual craniofacial asymmetry. If initiation is delayed until 3 to 6 months of age, the success rate of conservative management decreases to approximately 75\%, and the duration of therapy required to achieve resolution increases. For children who present after one year of age, or who fail conservative therapy, surgical release yields good-to-excellent outcomes in the majority of cases, although minor residual facial asymmetry may persist if skeletal maturity is reached before intervention.

For acute acquired torticollis (wry neck) due to minor musculoskeletal strain or position-related spasms, the prognosis is excellent. Most cases are self-limiting and resolve completely within 24 to 72 hours, with full recovery of cervical mobility. Refractory cases typically resolve within one to two weeks with conservative management (NSAIDs, heat therapy, and gentle stretching).

The prognosis for cervical dystonia (spasmodic torticollis) is more guarded. It is a chronic, progressive neurological condition that typically stabilizes after a few years of progression. Spontaneous, complete remission is rare, occurring in less than 10\% of patients, usually within the first 5 years of onset; these remissions are frequently temporary. However, with the advent of botulinum toxin therapy, the long-term outlook has improved significantly. Over 70\% to 80\% of patients experience significant symptomatic relief, reduction in pain, and improvement in neck posture and overall functional capability. Continuous management is required, but patients can maintain a good quality of life and remain active.

\section*{Prevention and Self-Care}
While congenital and neurological forms of torticollis cannot be directly prevented, simple preventative measures and self-care strategies can reduce the risk of acquired positional torticollis and aid in the recovery and management of existing conditions.

\begin{itemize}
    \item \textbf{Preventative Care for Infants:}
    \begin{itemize}
        \item \textit{Alternating Sleeping Positions:} Parents should place infants on their backs to sleep to reduce the risk of Sudden Infant Death Syndrome (SIDS). However, the direction of the infant's head should be alternated nightly (e.g., turning the head to the left one night, and to the right the next) to prevent positional plagiocephaly and muscle tightness.
        \item \textit{Environmental Adjustments:} Position toys, mobiles, and cribs so that the infant must turn their head in both directions to see interesting stimuli or interact with caregivers.
        \item \textit{Daily Tummy Time:} Implement supervised, daily prone positioning (tummy time) starting from birth. Begin with 1 to 2 minutes per session, building up to a total of 80 to 90 minutes per day by 3 months of age. This strengthens the cervical extensor muscles and promotes symmetrical motor development.
    \end{itemize}
    \item \textbf{Ergonomic and Self-Care Measures for Adults:}
    \begin{itemize}
        \item \textit{Workstation Ergonomics:} Arrange computer screens, chairs, and keyboards to maintain a neutral spine. The top of the monitor should be at eye level, directly in front of the user, to prevent chronic lateral tilting or rotation of the neck.
        \item \textit{Sleep Support:} Use an orthopedic or contour pillow that supports the natural curvature of the cervical spine. Avoid sleeping on the stomach, as this forces the neck into extreme rotation for prolonged periods.
        \item \textit{Regular Micro-breaks:} Individuals with sedentary desk jobs should stand, stretch, and perform gentle neck movements every 30 to 45 minutes to prevent muscle stiffness and positional strain.
        \item \textit{Stress Reduction:} Stress exacerbates muscle tension and can trigger or worsen spasms in patients with cervical dystonia. Practicing mindfulness, progressive muscle relaxation, yoga, or biofeedback can reduce systemic muscle tone and improve symptom control.
        \item \textit{Home Heat and Stretching:} Applying a warm compress or taking a warm shower relaxes hypertonic muscles, making gentle range-of-motion stretching exercises safer and more effective. Stretching should be slow and sustained; jerking or forcing the neck beyond comfort limits should be strictly avoided to prevent reflex spasms and tissue injury.
    \end{itemize}
\end{itemize}

\section*{Sources and Further Reading}
\begin{itemize}
    \item \textbf{Mayo Clinic} - Cervical Dystonia (Spasmodic Torticollis). Detailed patient care guides, diagnostic parameters, and treatment overviews. Available at: \texttt{https://www.mayoclinic.org/diseases-conditions/cervical-dystonia/symptoms-causes/syc-20354123}
    \item \textbf{National Institute of Neurological Disorders and Stroke (NINDS)} - Cervical Dystonia Information Page. Authoritative resources on the pathophysiology, clinical trials, and management of adult-onset spasmodic torticollis. Available at: \texttt{https://www.ninds.nih.gov/health-information/disorders/cervical-dystonia}
    \item \textbf{American Academy of Orthopaedic Surgeons (AAOS)} - Congenital Muscular Torticollis (Infant Wry Neck). OrthoInfo clinical library detailing infant stretch exercises, diagnosis, and surgical indications. Available at: \texttt{https://orthoinfo.aaos.org/en/diseases--conditions/congenital-muscular-torticollis-infant-wry-neck}
    \item \textbf{The Royal Children's Hospital Melbourne} - Torticollis. Comprehensive pediatric guide on positional torticollis, stretching exercises, and plagiocephaly prevention. Available at: \texttt{https://www.rch.org.au/kidsinfo/fact\_sheets/Torticollis/}
    \item \textbf{Dystonia Medical Research Foundation} - Cervical Dystonia Support and Research. Patient-focused advocacy group providing literature, peer support, and clinical guidelines on dystonic torticollis. Available at: \texttt{https://dystonia-foundation.org/what-is-dystonia/types-dystonia/cervical-dystonia/}
    \item \textbf{National Health Service (NHS) UK} - Wry Neck (Torticollis). General guidance on recognizing acute acquired wry neck, self-care strategies, and clinical management. Available at: \texttt{https://www.nhs.uk/conditions/wry-neck-torticollis/}
\end{itemize}